\section{Methods}

\subsection{Task Formulation}

We framed suggestion mining as binary sentence classification:
each sentence was labelled as suggestion~(1) or
non-suggestion~(0). Classification operated at the sentence
level, consistent with the SemEval-2019 task definition. We did
not use review-level context; each sentence was classified
independently.

\subsection{Tier 1: Regex Baseline}

The regex baseline classified a sentence as a suggestion if it
matched any of 24 hand-crafted regular expression patterns.
These patterns covered three categories:
modal suggestions (e.g.\ \texttt{\textbackslash bshould\textbackslash b},
\texttt{\textbackslash bcould improve\textbackslash b}, \texttt{\textbackslash
bneed to\textbackslash b}), imperatives (e.g.\ \texttt{\textbackslash
bplease\textbackslash b}, \texttt{\textbackslash brecommend\textbackslash b},
\texttt{\textbackslash bconsider\textbackslash b}), and conditionals (e.g.\
\texttt{\textbackslash bi wish\textbackslash b}, \texttt{\textbackslash bit
would be nice\textbackslash b}, \texttt{\textbackslash bhopefully\textbackslash
b}). Matching was case-insensitive. This baseline established a
performance floor and demonstrated the limitations of surface
patterns---for instance, ``We would recommend this hotel to
anyone'' matches the modal pattern but is not a suggestion.

\subsection{Tier 2: TF-IDF + LR Baseline}

The TF-IDF~+~LR baseline used a TF-IDF vectoriser (unigrams
and bigrams, maximum 10,000 features, minimum document
frequency~2, sublinear term-frequency scaling) followed by
logistic regression (regularisation $C=1.0$, balanced class
weights, maximum 1,000 iterations), implemented in scikit-learn
\citep{pedregosa2011scikit}. The model was trained on the SemEval training set. Balanced
class weights addressed the class imbalance by upweighting the
minority suggestion class. This baseline showed the value of
learned features over hand-crafted rules while remaining
interpretable and fast to train.

\subsection{Tier 3: BERT Classifier with Two-Stage Fine-Tuning}

Our primary model was \texttt{bert-base-uncased}
\citep{devlin2019bert}, a 110-million-parameter pre-trained
transformer implemented via the HuggingFace Transformers library
\citep{wolf2020transformers}, fine-tuned for binary
classification with an additional linear output layer.

\paragraph{Stage~1: SemEval fine-tuning.} We fine-tuned on the SemEval training
  set (8,500 sentences), using a 90/10 stratified train/validation split.
  Sentences were tokenised with a maximum length of 128 tokens. Training used the
  AdamW optimiser \citep{loshchilov2019adamw} with learning rate
  $2 \times 10^{-5}$, batch size~16,
  3~epochs, linear warmup over 10\% of training steps, and weight decay~0.01. We
  selected the checkpoint with the best macro~F1 on the validation split. On the
  SemEval validation split, BERT stage-1 achieved macro~F1 of 0.886 and
  suggestion-class~F1 of 0.828.

\paragraph{Stage~2: MBS domain adaptation.} Starting from the stage-1
  checkpoint, we further fine-tuned on the MBS train split, using an 85/15
  stratified train/validation split. The learning rate was reduced to $1 \times
  10^{-5}$ to avoid overwriting task knowledge learned in stage~1. All other
  hyperparameters remained the same. This two-stage approach followed the principle
  of progressive fine-tuning: pre-trained representations $\rightarrow$
  task-specific fine-tuning $\rightarrow$ domain-specific adaptation.

All training was performed on Apple Silicon (MPS backend) with a fixed random
seed of 42. Hyperparameters follow \citet{devlin2019bert}.

\subsection{Insight Generation: BERTopic Clustering}

To transform binary predictions into thematic insights, we
applied BERTopic \citep{grootendorst2022bertopic} to all
sentences predicted as suggestions by BERT stage-2. Sentence
embeddings were computed with \texttt{all-MiniLM-L6-v2}
\citep{reimers2019sentence}. Dimensionality reduction used UMAP
\citep{mcinnes2018umap} (5~components, 15~neighbours, cosine
metric, minimum distance~0.0). Clustering used HDBSCAN
\citep{campello2013hdbscan} (minimum cluster size~5). Topic
representations were extracted via a class-based TF-IDF
procedure. The number of topics was determined automatically.

As a complementary analysis, we applied keyword-based aspect
grouping to the same set of predicted suggestions. Six
hotel-relevant aspects were defined: room, food and beverage,
service, facilities, location, and value. Each sentence was
assigned to the aspect whose keywords appeared most frequently,
or to ``other'' if no keywords matched. This provided a coarser
but more interpretable categorisation than BERTopic.

\subsection{Evaluation}

Our primary metric is suggestion-class~F1, computed on label~1 only. We also
report precision, recall, and macro~F1. All models are evaluated on both the
SemEval test set (in-domain evaluation) and the MBS test split (cross-domain
evaluation). The difference between in-domain and cross-domain
suggestion-class~F1---the performance delta---quantifies the domain gap.
We do not perform a separate manual validation step on the full
MBS predictions. This limitation is discussed in
Section~6.3.
